\documentclass[11pt]{letter}

\usepackage[margin=0.5in]{geometry}
\usepackage{hyperref}

\begin{document}

\begin{letter}{}

\opening{Dear Hiring Team,}

I am writing to apply for the Machine Learning Engineer - Enterprise position at Boson AI in Toronto. I am an M.Sc. Computer Science graduate from the University of Calgary with experience building Python/C++ systems for machine learning, scalable data processing, backend services, and AI workflow automation. Boson AI's focus on enterprise AI, LLM-powered software, model evaluation, search systems, and agentic workflows is closely aligned with the technical direction I have been building toward through my graduate research and recent automation projects.

My strongest machine learning experience comes from research on digital elevation model super-resolution. I built Python and PyTorch pipelines for data collection, preprocessing, model training, and evaluation; trained deep ResNet models for 4x terrain reconstruction; and evaluated quality using RMSE, slope, TRI, TPI, and wavelet-based metrics. This work required careful experimentation, data-quality checks, and attention to whether model results generalized beyond convenient random splits. That evaluation mindset is one I would bring to Boson AI's work on benchmarking models, writing evals, and improving reliability in enterprise-facing AI systems.

I also bring systems experience from my graduate work with BigGeo and Mitacs, where I developed C++/Python software for scalable processing, storage, and retrieval of Earth observation data. I designed tensor-based storage and retrieval methods, implemented multiresolution data workflows, and improved processing performance through C++ multithreading. This background gives me a practical appreciation for building software that is not only technically interesting, but also efficient, reliable, and usable at scale.

More recently, I have been building an agentic job-search automation platform using OpenClaw, n8n, LLM workflows, Docker, SearXNG, Google Sheets, Google Drive, and Telegram integrations. The system searches for jobs, analyzes fit, tracks opportunities, and prepares tailored application materials. While this is not the same as deploying large enterprise AI products, it has given me hands-on exposure to orchestrating LLM-based workflows, connecting models with search and external tools, and designing automation that has to remain auditable and useful rather than merely impressive.

Before graduate school, I worked as a back-end developer building Django REST APIs for NLP, text-to-speech, and speech-to-text services backed by PostgreSQL. I also contributed to a WebRTC-based chatbot system using GPT-style responses and FAQ data. That experience complements my research background with practical software engineering habits: API design, authentication flows, database-backed features, integration work, and debugging across a production-oriented backend.

I am excited by this role because it combines machine learning depth with the engineering discipline needed to deliver enterprise AI products. I would be glad to contribute my PyTorch, Python, backend, evaluation, and AI workflow experience while continuing to grow in areas such as retrieval-augmented generation, model orchestration, scalable inference, and enterprise integration.

Thank you for considering my application. I would welcome the opportunity to discuss how my background in machine learning, scalable systems, and AI workflow automation can contribute to Boson AI.

\closing{Sincerely,

Amir Mirzai Golpayegani

\href{mailto:amirgolpaa24@gmail.com}{amirgolpaa24@gmail.com}}

\end{letter}

\end{document}
